\question{} 

เรากำหนดให้มีกล่องจำนวนหนึ่งวางอยู่บนโต๊ะ ซึ่งกล่องบางใบอาจจะถูกซ้อนอยู่ภายในกล่องใบอื่นก็ได้ และมีเงื่อนไขเพื่อความชัดเจนว่า
\begin{itemize}
    \item \uline{ถ้า}กล่อง X ถูกซ้อนในกล่อง Y \uline{และ}กล่อง Y ถูกซ้อนในกล่อง Z 
        \uline{แล้ว}กล่อง X จะถือว่าถูกซ้อนในกล่อง Z ด้วย
    \item \uline{ถ้า}กล่อง X ถูกซ้อนในกล่อง Y \uline{แล้ว}กล่อง Y จะไม่ถูกซ้อนในกล่อง X
\end{itemize}
นอกจากเงื่อนไขข้างต้นนี้ เรากำหนดนิยามเพิ่มเติมดังต่อไปนี้
\begin{quote}
    \textbf{\uline{นิยาม}}\; หากกำหนดให้ $B$ เป็นสัญลักษณ์ที่สื่อถึงรูปแบบของการจัดวางกล่องบนโต๊ะตามเงื่อนไขที่กำหนดให้ 
    แล้วสัญลักษณ์ $B[n]$ จะหมายถึง ``จำนวนของกล่องทั้งหมดบนโต๊ะที่มีกล่องอื่น ๆ จำนวน $n$ ใบซ้อนอยู่ภายใน''

    \marginnote{%
        \centering
        \includegraphics[width=0.98\linewidth]{figures/quickfire_central_regional_nestbox.eps}\\
        รูปตัวอย่างการจัดวางกล่อง $B_0$
    }

    ยกตัวอย่างเช่น สมมติว่า $B_0$ คือการจัดวางกล่อง 5 ใบบนโต๊ะรูปแบบหนึ่ง ดังที่แสดงทางด้าน\ifpageodd{ขวา}{ซ้าย}มือนี้
    เราสามารถบอกได้ว่า $B_0[1] = \mathrm{2}$ (กล่องสีฟ้า), $B_0[2] = \mathrm{0}$, 
    $B_0[3] = \mathrm{0}$ และ $B_0[4] = \mathrm{1}$ (กล่องสีแดง) ฯลฯ
\end{quote}

\noindent
\textbf{\uline{โจทย์}}\; จงคำนวณว่าสำหรับการจัดวางกล่อง $B_1$ \uline{ในรูปแบบใด\,ๆ} 
ที่สอดคล้องกับเงื่อนไขว่า $B_1[1] = \mathrm{5}$,\, $B_1[2] = \mathrm{4}$,\, $B_1[3] = \mathrm{3}$,\, 
$B_1[4] = \mathrm{2}$ และ $B_1[5] = \mathrm{1}$ จะต้อง\uline{มีจำนวนกล่องรวม}ทั้งหมดอย่างน้อยกี่ใบ?
